\label{sec:chain}
\subsection{What the dictionary says about the chain}
\begin{enumerate}[leftmargin=*]
\item \textbf{The dilation must be carried, because it cannot be pinned.}
\begin{itemize}
\item In a bias-free network the dilation is a free direction of every contact configuration (Proposition~\ref{prop:dilfree}) and the
modular flow of the non-commutative strata (Theorem~\ref{thm:modular}).
\item So no closure built from local (contact) data fixes it, and errors along it are not damped. That is note XLIV's critical mode,
now seen from three sides.
\item The remedy is stage 13's: carry the charge (the radial partition function's second cumulant) and write the variance, $\kappa_3$ and
$\kappa_4$ along the dilation as one package.
\item \emph{Prediction:} networks with biases show a damped gain mode.
\end{itemize}
\item \textbf{Resolution is set by $r^2/4$, so effort should follow the wall density.} By Corollary~\ref{cor:resell}, a neuron at field $r$ needs chaos
order about $r^2/4$ before its kink is seen at all, and contributes at most $\varphi(r)/r^3$.
\begin{itemize}
\item The chain's truncation error at order $K$ concentrates on neurons with $r^2\lesssim4K$, weighted by $\varphi(r)K^{-3/2}$.
\item \emph{Prediction (testable on note XLI's dumps):} the per-neuron benefit of v56's chaos correction is supported on $|r|\lesssim2\sqrt3$,
with a $\varphi(r)$-weighted profile.
\item \emph{Prediction:} the next chaos order removes about $35\%$ of the remaining chaos-truncation energy.
\end{itemize}
\item \textbf{Positivity is a sticky constraint.} The chain's state is a functional on the truncated operator system $\mathcal E_K$. If it
leaves the positive cone, the pseudo-law has negative mass somewhere near a wall.
\begin{itemize}
\item Klartag's rule is to project each layer's update onto the face it touches and continue there.
\item This is truth-free and costs $O(n^2)$ per layer for the marginal (one-neuron) moment matrices.
\item A first diagnostic: count how often v56's per-neuron moment matrices (orders $\le4$) fail to be positive semidefinite.
\end{itemize}
\item \textbf{The cell-adapted Gaussian.} The maximal resolved probe at an input is the natural local Gaussian of its cell, and its frame
code records the cell's local wall geometry. It is a candidate local model for future hybrid estimators. Sampling is roughly
$1000\times$ too expensive at this budget (stage 13), so this is no algorithm yet.
\end{enumerate}

\subsection{Open problems}
\begin{enumerate}[leftmargin=*]
\item \textbf{Klartag's estimate for walls.} Prove that, for Gaussian rows, the sticky localization reaches a resolved probe of
log-volume $\ge c\,\cdot$(free dimensions) before contacts saturate. Klartag's $n^2$-versus-$n$ gain would be the ratio between the
best Gaussian probe and the best ball inside a typical cell.
\item \textbf{Twisted gap labels.} Compute the dilation-twisted (KMS) pairing of the bent strata, the non-commutative part of the
gap-labelling group, and decide whether it equals a correction term of the mean.
\item \textbf{Left and right.} Make precise whether the Tomita conjugation $JL(G)J=R(G)$ at bent strata exchanges stage 11's forward and
backward harmonic states.
\item \textbf{The radial partition function.} Find a recursion for $Z_l(\beta)$ through a layer. The dilation package of stage 13 uses
only its second cumulant.
\item \textbf{Codes across inputs.} The tight frames of Proposition~\ref{prop:john} vary with the input. Is their distribution over the
sphere a function of the internal clock alone, as stage 12's field law is?
\end{enumerate}
